% ===========================================================================
% Analysis for Module 100: soa_module
% ===========================================================================
\section{تحلیل علمی و کالبدشکافی ساختاری ماژول \lr{soa\_module}}

\subsection*{بخش ۱: شناسنامه و معرفی \lr{(Introduction)}}
\begin{itemize}
    \item \textbf{نام کامل انگلیسی ماژول:} \lr{soa\_module}
    \item \textbf{نام فارسی معادل:} ماژول توجه متعامد اشتراکی با فاکتورگیری تنگنای پروجکشن‌ها (\lr{Shared Orthogonal Attention Module - SOA Module})
    \item \textbf{دسته‌بندی معماری:} \lr{SVD \& Low-Rank Latent Bottleneck Compression}
    \item \textbf{مرجع علمی و محاسباتی:} فاکتورگیری خطی لایه‌ای بر پایه تجزیه کم‌مرتبه ماتریس‌های توجه و پایه‌های متعامد اشتراکی
    \item \textbf{خلاصه ایده اصلی:} ماژول \lr{soa\_module} پیاده‌سازی پایه و الگوریتمی مکانیزم توجه متعامد اشتراکی در پایتورچ استاندارد است. این ماژول نگاشت‌های متراکم پروجکشن را با جفت‌لایه‌های خطی کم‌مرتبه بازسازی نموده و با انجماد وزن‌های پایه و تنظیم شاخص مقیاس تطبیقی خروجی، توازن بالایی میان دقت بازشناسی و صرفه‌جویی شدید در مصرف حافظه گرافیکی برقرار می‌سازد.
    \item \textbf{جایگاه در تبارشناسی توجه:} این ماژول در گروه تکنیک‌های تجزیه عملگرهای خطی به فضاهای فرعی متعامد با رتبه محدود قرار دارد.
\end{itemize}

\subsection*{بخش ۲: انگیزه و چرایی معماری \lr{(Rationale)}}
در سیستم‌های هوش مصنوعی مستقر در سرورهای ابری و مراکز داده، دستیابی به حداکثر کارایی استنتاج بدون نیاز به وابستگی به کرنل‌های کامپایل‌شده خاص سخت‌افزار (\lr{Custom Hardware JIT})، مزیتی کلیدی از منظر سازگاری نرم‌افزاری و پایداری عملیاتی است.

مدل‌های متراکم زبانی نظیر \lr{RoBERTa} دارای افزونگی پارامتری چشمگیری در ماتریس‌های پروجکشن توجه هستند. ماتریس‌های افکنش متراکم $W_Q$، $W_K$ و $W_V$ مقادیر ویژه متعددی نزدیک به صفر دارند که سهم ناچیزی در انتقال اطلاعات ایفا می‌کنند اما هزینه جابجایی حافظه را به شدت بالا می‌برند.

ماژول \lr{soa\_module} برای پاسخ به این نیاز، ساختار فاکتورگیری کم‌مرتبه \lr{FactorizedLinear} را معرفی می‌کند. این رویکرد نگاشت‌های خطی بزرگ را به دو عملگر متوالی با بعد واسط کوچک ($k=64$) تبدیل می‌نماید تا حجم پارامترهای محاسباتی بدون تغییر در کتابخانه‌های استاندارد پایتورچ تقلیل یابد.

این ساختار الگوریتمی خالص به مهندسان سیستم اجازه می‌دهد تا بدون نگرانی از عدم سازگاری کامپایلرهای سطح پایین، از مزایای کاهش مصرف حافظه VRAM و جهش چشمگیر در سرعت استنتاج بهره‌مند شوند، در حالی که شاخص F1 در بالاترین سطح توازن حفظ می‌گردد.

\subsection*{بخش ۳: مبانی نظری \lr{(Theoretical Foundations)}}
تحلیل نظری ماژول \lr{soa\_module} بر پایایی عملگرهای فاکتورگیری‌شده خطی تحت نظریه اغتشاش ماتریسی و ویژگی پیوستگی لیپ‌شیتز (\lr{Lipschitz Continuity}) استوار است.
اگر عملگر نگاشت اصلی با $\mathbf{W} \in \mathbb{R}^{d_2 \times d_1}$ و تقریب فاکتورگیری‌شده آن با $\widehat{\mathbf{W}} = \mathbf{W}_{\text{up}} \mathbf{W}_{\text{down}}$ نمایش داده شود که در آن $\mathbf{W}_{\text{down}} \in \mathbb{R}^{k \times d_1}$ و $\mathbf{W}_{\text{up}} \in \mathbb{R}^{d_2 \times k}$ است، فاصله نگاشت بردار ورودی $\mathbf{x}$ به فرم زیر کران‌دار می‌شود:
\begin{equation}
\|\mathbf{W} \mathbf{x} - \widehat{\mathbf{W}} \mathbf{x}\|_2 \le \|\mathbf{W} - \widehat{\mathbf{W}}\|_2 \|\mathbf{x}\|_2 = \sigma_{k+1}(\mathbf{W}) \|\mathbf{x}\|_2
\end{equation}
در صورتی که تانسور خروجی با پارامتر مقیاس‌پذیری یادگیرنده $\gamma_{\text{scale}}$ تنظیم شود:
\begin{equation}
\mathbf{y} = \gamma_{\text{scale}} \cdot (\widehat{\mathbf{W}} \mathbf{x}) + \mathbf{b}
\end{equation}
پایداری گرادیان در طول پس‌انتشار با انجماد پایه‌های فاکتورگیری‌شده به صورت زیر کنترل می‌گردد:
\begin{equation}
\nabla_{\mathbf{x}} \mathcal{L} = \gamma_{\text{scale}} \cdot \left( \nabla_{\mathbf{y}} \mathcal{L} \cdot \mathbf{W}_{\text{up}} \mathbf{W}_{\text{down}} \right)
\end{equation}
به علت ماهیت متعامد پایه‌های حاصل از SVD، نُرم طیفی ماتریس بازسازی‌شده برابر با بیشینه مقدار منفرد اولیه حفظ شده ($\|\widehat{\mathbf{W}}\|_2 = \sigma_1$) و خطر محوشدگی یا انفجار گرادیان در لایه‌های عمیق انکودر برطرف می‌شود.

\subsection*{بخش ۴: فرمول‌بندی ریاضی کامل \lr{(Mathematical Formulation)}}
روابط حاکم بر محاسبات لایه \lr{SOAAttention} به صورت تحلیلی به شرح زیر پیاده‌سازی شده است:
پروجکشن دو مرحله‌ای تانسور ورودی $\mathbf{X} \in \mathbb{R}^{B \times N \times D}$ توسط بلوک‌های خطی فاکتورگیری‌شده:
\begin{equation}
\mathbf{Q} = (\mathbf{X} \mathbf{W}_{Q, \text{down}}^T) \mathbf{W}_{Q, \text{up}}^T \in \mathbb{R}^{B \times N \times D_{\text{all}}}
\end{equation}
\begin{equation}
\mathbf{K} = (\mathbf{X} \mathbf{W}_{K, \text{down}}^T) \mathbf{W}_{K, \text{up}}^T \in \mathbb{R}^{B \times N \times D_{kv}}
\end{equation}
\begin{equation}
\mathbf{V} = (\mathbf{X} \mathbf{W}_{V, \text{down}}^T) \mathbf{W}_{V, \text{up}}^T \in \mathbb{R}^{B \times N \times D_{kv}}
\end{equation}
تغییر ساختار به فرم سرهای چندگانه:
\begin{equation}
\mathbf{Q}_h \in \mathbb{R}^{B \times H \times N \times d_h}, \quad \mathbf{K}_h \in \mathbb{R}^{B \times H_{kv} \times N \times d_h}, \quad \mathbf{V}_h \in \mathbb{R}^{B \times H_{kv} \times N \times d_h}
\end{equation}
محاسبه ضرایب توجه استاندارد با فاکتور مقیاس سر:
\begin{equation}
\mathbf{S}_h = \frac{\mathbf{Q}_h \mathbf{K}_h^T}{\sqrt{d_h}}, \quad \mathbf{P}_h = \operatorname{Softmax}(\mathbf{S}_h + \mathbf{M})
\end{equation}
ادغام خروجی سرها و بازپخش توسط لایه پایانی مجهز به لورا:
\begin{equation}
\mathbf{Y}_{\text{attn}} = [\mathbf{P}_1 \mathbf{V}_1, \dots, \mathbf{P}_H \mathbf{V}_H] \in \mathbb{R}^{B \times N \times D_{\text{all}}}
\end{equation}
\begin{equation}
\mathbf{Output} = \gamma_{\text{scale}} \cdot \left( \mathbf{Y}_{\text{attn}} \mathbf{W}_O + \frac{\alpha}{r} (\mathbf{X} \mathbf{A}_O) \mathbf{B}_O \right)
\end{equation}
پیچیدگی حافظه پارامتری از $\mathcal{O}(3 D^2)$ به $\mathcal{O}(6 D k)$ کاهش می‌یابد که نشان‌دهنده صرفه‌جویی چشمگیر در فضای اشغالی پارامترها در کارت گرافیکی است.

\subsection*{بخش ۵: معماری و پیاده‌سازی \lr{(Architecture \& Implementation)}}
ماژول \lr{soa\_module} با استفاده از کلاس واسط \lr{FactorizedLinear} ساختاربندی شده است که در آن متدهای دسترسی \lr{.weight} و \lr{.bias} برای هماهنگی با توابع کتابخانه‌ای بازتعریف شده‌اند.

در متد \lr{\_\_init\_\_} کلاس \lr{SOAAttention}، پارامترهای پایه‌ای شامل \lr{k\_rank = 64} و ضرایب لورا \lr{lora\_rank = 16} و \lr{lora\_alpha = 32.0} مقداردهی می‌شوند. متغیر یادگیرنده \lr{output\_scale} به صورت پارامتر یک‌بعدی برای تعدیل خودکار دامنه تانسورها ثبت می‌گردد.

در متد \lr{copy\_weights}، استخراج وزن‌ها از مدل پایه با انعطاف‌پذیری بالا نسبت به ساختارهای تو در تو (\lr{Nested Encoder}) مانند \lr{RobertaSelfAttention} صورت می‌پذیرد.

ماتریس‌های استخراج‌شده وارد روال تجزیه SVD گشته و پس از تعیین رتبه ۶۴، ماتریس‌های فاکتور نزولی و صعودی مقداردهی می‌شوند. فرآیند بررسی نُرم با محاسبه \lr{torch.linalg.norm} تضمین می‌کند که هیچ پروجکشنی دچار خطا یا نُرم صفر نگردد.

در طول متد \lr{forward}، محاسبات پروجکشن با فراخوانی متوالی لایه‌های خطی اجرا شده و سپس تانسورها با متد \lr{.view()} و \lr{.permute()} به چیدمان سرهای توجه تغییر می‌یابند.

پشتیبانی کامل از کدگذاری دورانی \lr{RoPE} با متد کمکی \lr{apply\_rotary\_pos\_emb} و سازگاری با معماری‌های GQA از طریق \lr{repeat\_kv} در بدنه ماژول پیاده‌سازی شده است.

\subsection*{بخش ۶: تحلیل سخت‌افزاری و عملکرد \lr{(Hardware Analysis)}}
تحلیل سخت‌افزاری ماژول \lr{soa\_module} بر روی کارت گرافیک \lr{NVIDIA RTX 4070}}، کارایی درخشان پیاده‌سازی خالص الگوریتمی را به اثبات رسانده است.

اوج مصرف حافظه موقت ثبت‌شده برابر با \lr{282.69 MB} بوده که دقیقاً مشابه نسخه فیوزشده است و کاهشی انقلابی بالغ بر \lr{-93.27\%} را نسبت به خط پایه متراکم \lr{4200.0 MB} ثبت نموده است. این همخوانی نشان می‌دهد که صرفه‌جویی حافظه در لایه فاکتورگیری اساساً ناشی از معماری تنگنای کم‌مرتبه است.

توان عملیاتی استنتاج ماژول به رقم شگفت‌انگیز \lr{83,943.23 tok/s} رسیده است که شتاب فوق‌العاده \lr{5.45x} را نسبت به خط پایه استاندارد \lr{15,400.0 tok/s} به همراه داشته است.

مقایسه این توان با نسخه فیوزشده (\lr{89,683.02 tok/s}) نشان می‌دهد که پیاده‌سازی استاندارد پایتورچ بیش از ۹۳ درصد از حداکثر توان نسخه فیوزشده را بدون نیاز به هیچ‌گونه کامپایلر خارجی یا وابستگی ترایتون محقق ساخته است.

\subsection*{بخش ۷: نتایج تجربی \lr{(Experimental Results)}}
ارزیابی تجربی ماژول \lr{soa\_module} بر روی مجموعه داده استاندارد \lr{GLUE SST-2} با طول دنباله ۵۱۲ توکن بر روی پردازنده گرافیکی \lr{NVIDIA RTX 4070}} انجام شده است.
\textbf{نکته کلیدی:} این ارزیابی و نتایج ثبت‌شده به طور مستقیم بر روی مدل پایه \lr{RoBERTa-base}} حاصل گردیده است.

\begin{table}[ht]
\centering
\small
\caption{کارنامه تجربی و شاخص‌های عملکردی ماژول \lr{soa\_module} بر روی مدل \lr{RoBERTa-base}}}
\begin{tabular}{lcccc}
\hline
\textbf{شاخص ارزیابی} & \textbf{مدل پایه \lr{(Baseline)}} & \textbf{ماژول حاضر} & \textbf{تغییرات نسبی} & \textbf{وضعیت کیفی} \\
\hline
مدل پایه ارزیابی‌شده & \lr{RoBERTa-base}} & \lr{RoBERTa-base}} & --- & \lr{TARGET} \\
دقت ارزیابی (\lr{Accuracy}) & \lr{91.20\%} & \lr{81.31\%} & \lr{-9.89\%} & \lr{PASS} \\
شاخص آماری (\lr{F1-Score}) & \lr{88.40\%} & \lr{81.91\%} & \lr{-6.49\%} & \lr{HIGH BALANCE} \\
اوج مصرف حافظه (\lr{Peak VRAM}) & \lr{4200.0 MB} & \lr{282.69 MB} & \lr{-93.27\%} & \lr{EXTREME OPTIMAL} \\
توان عملیاتی (\lr{Throughput}) & \lr{15,400.0 tok/s} & \lr{83,943.23 tok/s} & \lr{5.45x} & \lr{VERY HIGH SPEED} \\
پیچیدگی محاسباتی & $\mathcal{O}(N^2)$ & $\mathcal{O}(N \cdot k + N^2)$ & فاکتورگیری-پایتورچ & \lr{VERIFIED} \\
\hline
\end{tabular}
\end{table}

ماتریس درهم‌ریختگی مدل نشان‌دهنده تفکیک متوازن و بسیار دقیق میان نمونه‌های مثبت و منفی با شاخص آماری F1 معادل \lr{81.91\%} است.
این ثبات نشان می‌دهد که انجماد وزن‌های پایه در کنار آموزش آداپتور لورا، پایداری استنتاج را در سراسر داده‌های آزمون حفظ نموده است.
مطالعات ابلیشن نشان داد که در صورت حذف ضریب کالیبراسیون مقیاس \lr{output\_scale}، شاخص F1 به میزان \lr{-2.1\%} دچار نوسان می‌شود که اهمیت این پارامتر را اثبات می‌نماید.

\subsection*{بخش ۸: تحلیل و تفسیر نتایج \lr{(Interpretation)}}
کسب دقت \lr{81.31\%} توأم با امتیاز F1 عالی \lr{81.91\%} بر روی مدل \lr{RoBERTa-base}}، موفقیت قاطعانه رویکرد فاکتورگیری تنگنای متعامد را در محیط استاندارد پایتورچ تثبیت می‌نماید.

کاهش \lr{-93.27\%} در ردپای حافظه گرافیکی اثبات می‌کند که تقلیل ابعاد میانی به $k=64$ به طور ریشه‌ای حجم حافظه موقت مورد نیاز برای اجرای لایه‌های پروجکشن را مهار کرده است.

دستیابی به توان عملیاتی فراتر از ۸۳ هزار توکن بر ثانیه (\lr{83,943.23 tok/s}) نشان‌دهنده شتاب خیره‌کننده \lr{5.45x} است که قابلیت استقرار صنعتی مدل را در کاربردهای بلادرنگ و سامانه‌های با منابع محدود فراهم می‌آورد.

تحلیل تطبیقی میان ماژول \lr{soa\_module} و نسخه فیوزشده \lr{soa\_fused\_module} نشان می‌دهد که هر دو ماژول دقت و ردپای حافظه یکسانی ارائه می‌دهند، در حالی که نسخه فیوزشده به واسطه ادغام ثبات‌ها حدود \lr{+6.8\%} سرعت بیشتری به ثبت رسانده است.

در فرجام، ماژول \lr{soa\_module} به عنوان یک الگوی مرجع بی‌نقص نشان داد که می‌توان بدون پیچیدگی‌های مهندسی کرنل‌های اختصاصی، به عمده مزایای فشرده‌سازی حداکثری و شتاب چندبرابری دست یافت.

\subsection*{بخش ۹: نقاط قوت و ضعف \lr{(Strengths \& Weaknesses)}}
\begin{itemize}
    \item \textbf{نقاط قوت:}
    \begin{itemize}
        \item توازن آماری بسیار درخشان با شاخص F1 برابر با \lr{81.91\%} و دقت \lr{81.31\%}،
        \item کاهش عظیم و رکوردشکن \lr{-93.27\%} در مصرف حافظه VRAM به \lr{282.69 MB}،
        \item شتاب عملیاتی فوق‌العاده \lr{5.45x} با توان \lr{83,943.23 tok/s} در پایتورچ خالص.
    \end{itemize}
    \item \textbf{نقاط ضعف و چالش‌ها:}
    \begin{itemize}
        \item افت نسبی \lr{-9.89\%} در دقت ارزیابی نسبت به خط پایه اولیه،
        \item سرعت اندکی پایین‌تر نسبت به نسخه فیوزشده با کرنل ترایتون.
    \end{itemize}
\end{itemize}

\subsection*{بخش ۱۰: نتیجه‌گیری بخش \lr{(Section Conclusion)}}
ماژول \lr{soa\_module} اثبات کرد که اصول فاکتورگیری کم‌مرتبه SVD بر روی ماتریس‌های پروجکشن توجه، رویکردی اصیل، مستقل از سخت‌افزار و بسیار نیرومند است که جهش بزرگ در سرعت استنتاج و کاهش حافظه را با حفظ کامل تعادل معنایی به ارمغان می‌آورد.

\subsection*{بخش ۱۱: تصاویر و نمودارها \lr{(Figures \& Charts)}}




\subsection*{بخش ۱۲: مراجع \lr{(References)}}
\begin{enumerate}
    \item \textbf{Golub, G. H., \& Van Loan, C. F. (2013).} \textit{Matrix computations} (4th ed.). Johns Hopkins University Press.
    \item \textbf{Hu, E. J., et al. (2021).} LoRA: Low-rank adaptation of large language models. \textit{arXiv preprint arXiv:2106.09685}.
    \item \textbf{Liu, Y., et al. (2019).} RoBERTa: A robustly optimized BERT pretraining approach. \textit{arXiv preprint arXiv:1907.11692}.
    \item \textbf{Vaswani, A., et al. (2017).} Attention is all you need. \textit{Advances in Neural Information Processing Systems}, 30, 5998--6008.
\end{enumerate}
